% Original source: Zhou, physical PDF pp. 173--177 (printed pp. 157--161).
% Figures: none. Tables: Table 6.3 is native LaTeX.

% ZHOU-SOURCE-PAGE: 173 PRINTED: 157

smile for the Black-Scholes pricing model?

\solution
Implied volatility is the volatility that makes the model option price equal to the
market option price. Volatility smile describes the relationship between the implied
volatility of the options and the strike prices for a given asset. For currency options,
implied volatilities tend to be higher for in-the-money and out-of-the-money options
than for at-the-money options. For equity, volatility often decreases as the strike price
increases (also called volatility skew). The Black-Scholes model assumes that the asset
price follows a lognormal distribution with constant volatility. In reality, volatilities are
neither constant nor deterministic. In fact, the volatility is a stochastic process itself.
Furthermore, there may be jumps in asset prices.

\emph{B.} You have to price a European call option either with a constant volatility 30\%
or by drawing volatility from a random distribution with a mean of 30\%. Which option
would be more expensive?

\solution
Many would simply argue that stochastic volatility makes the stock price more volatile,
so the call price is more valuable when the volatility is drawn from a random distribution.
Mathematically, the underlying argument is that the price of a European call option is a
convex function of volatility and as a result

\[
c(E[\sigma])\leq E[c(\sigma)],
\]

where $\sigma$ is the random variable representing volatility and $c$ is the call option
price. Is the underlying argument correct? It's correct in most, but not all, cases. If the
call price $c$ is always a convex function of $\sigma$, then

\[
\frac{\partial^2c}{\partial\sigma^2}\geq0.
\]

$\frac{\partial c}{\partial\sigma}$ is the Vega of the option. For a European call option,

\[
\nu=\frac{\partial c}{\partial\sigma}=S\sqrt{\tau}N'(d_1)
=\frac{S\sqrt{\tau}}{\sqrt{2\pi}}\exp(-d_1^2/2).
\]

The secondary partial derivative $\frac{\partial^2c}{\partial\sigma^2}$ is called Volga.
For a European call option,

\[
\frac{\partial^2c}{\partial\sigma^2}
=\frac{S\sqrt{\tau}}{\sqrt{2\pi}}\exp(-d_1^2/2)\frac{d_1d_2}{\sigma}
=\nu\frac{d_1d_2}{\sigma}.
\]

$\nu$ is always positive. For most out-of-the-money call options, both $d_1$ and $d_2$
are negative; for most in-the-money call options, both $d_1$ and $d_2$ are positive. So
$d_1d_2>0$ in most cases and $c$ is a convex function of $\sigma$ when $d_1d_2>0$.
But theoretically, we can have conditions that $d_1>0$ and $d_2<0$ and
$\frac{\partial^2c}{\partial\sigma^2}<0$ when the option is close to being

% ZHOU-SOURCE-PAGE: 174 PRINTED: 158

at-the-money. So the function is not always convex. In those cases, the option with
constant volatility may have a higher value.

\emph{C.} The Black-Scholes formula for non-dividend paying stocks assumes that the stock
follows a geometric Brownian motion. Now assume that you don't know the stochastic
process followed by the stock price, but you have European call prices for all (continuous)
strike prices $K$. Can you determine the risk-neutral probability density function of the
stock price at time $T$?

\solution
The payoff of a European call at its maturity date is $\operatorname{Max}(S_T-K,0)$.
Therefore under risk-neutral measure, we have

\[
c=e^{-r\tau}\int_K^\infty(s-K)f_{S_T}(s)\,ds,
\]

where $f_{S_T}(s)$ is the probability density function of $S_T$ under the risk-neutral
probability measure. Taking the first and second derivatives of $c$ with respect to $K$,
\footnote[10]{To calculate the derivatives requires the Leibniz integral rule, a formula for differentiating a definite integral whose limits are functions of the differential variable:
\(\displaystyle \frac{\partial}{\partial z}\int_{a(z)}^{b(z)}f(x,z)\,dx
=\int_{a(z)}^{b(z)}\frac{\partial f(x,z)}{\partial z}\,dx
+f(b(z),z)\frac{\partial b}{\partial z}-f(a(z),z)\frac{\partial a}{\partial z}.\)} we have

\[
\frac{\partial c}{\partial K}
=e^{-r\tau}\frac{\partial}{\partial K}\int_K^\infty(s-K)f_{S_T}(s)\,ds
\]

\[
=e^{-r\tau}\int_K^\infty\frac{\partial(s-K)}{\partial K}f_{S_T}(s)\,ds
-e^{-r\tau}(K-K)\times1
=e^{-r\tau}\int_K^\infty-f_{S_T}(s)\,ds
\]

and

\[
\frac{\partial^2c}{\partial K^2}
=\frac{\partial}{\partial K}\left(\frac{\partial c}{\partial K}\right)
=e^{-r\tau}\frac{\partial}{\partial K}\int_K^\infty-f_{S_T}(s)\,ds
=e^{-r\tau}f_{S_T}(K).
\]

Hence the risk-neutral probability density function is

\[
f_{S_T}(K)=e^{r\tau}\frac{\partial^2c}{\partial K^2}.
\]

\subsection{Option Portfolios and Exotic Options}

In addition to the pricing and properties of vanilla European and American options, you
may be expected to be familiar with the construction and payoff of basic option-based
trading strategies---covered call, protective put, bull/bear spread, butterfly spread,
straddle, etc. Furthermore, if you are applying for a derivatives-related position, you

% ZHOU-SOURCE-PAGE: 175 PRINTED: 159

should also have a good understanding of pricing and hedging of some of the common
exotic derivatives---binary option, barrier option, Asian option, chooser option, etc.

\subsubsection*{Bull spread}

What are the price boundaries for a bull call spread?

\solution
A bull call spread is a portfolio with two options: long a call $c_1$ with strike $K_1$
and short a call $c_2$ with strike $K_2$ ($K_1<K_2$). The cash flow of a bull spread is
summarized in table 6.3.

\begin{table}[H]
\centering
\small
\begin{tabular}{|l|c|c|c|c|}
\hline
Cash flow & Time 0 & \multicolumn{3}{c|}{Maturity $T$} \\
\cline{3-5}
 & & $S_T\leq K_1$ & $K_1<S_T<K_2$ & $S_T\geq K_2$ \\
\hline
Long $c_1$ & $-c_1$ & 0 & $S_T-K_1$ & $S_T-K_1$ \\
\hline
Short $c_2$ & $c_2$ & 0 & 0 & $-(S_T-K_2)$ \\
\hline
Total & $c_2-c_1<0$ & 0 & $S_T-K_1$ & $K_2-K_1$ \\
\hline
\end{tabular}
\caption{Cash flows of a bull call spread.}
\label{tab:zhou-table-6-3}
\end{table}

Since $K_1<K_2$, the initial cash flow is negative. Considering that the final payoff is
bounded by $K_2-K_1$, the price of the spread, $c_1-c_2$, is bounded by
$e^{-rT}(K_2-K_1)$. Besides, the payoff is also bounded by
$\frac{K_2-K_1}{K_2}S_T$, so the price is also bounded by
\[
\frac{K_2-K_1}{K_2}S.
\]

\subsubsection*{Straddle}

Explain what a straddle is and when you want to purchase a straddle.

\solution
A straddle includes long positions in both a call option and a put option with the same
strike price $K$ and maturity date $T$ on the same stock. The payoff of a long straddle is
$|S_T-K|$. So a straddle may be used to bet on large stock price moves. In practice, a
straddle is also used as a trading strategy for making bets on volatility. If an investor
believes that the realized (future) volatility should be much higher than the implied
volatility of call and put options, he or she will purchase a straddle. For example,

% ZHOU-SOURCE-PAGE: 176 PRINTED: 160

the value of an at-the-money call or put is almost a linear function of volatility. If the
investor purchases an at-the-money straddle, both the call and the put options have the
price $c\approx p\approx0.4\sigma_iS\sqrt{\tau}$, where $\sigma_i$ is the implied
volatility. If the realized volatility $\sigma_r>\sigma_i$, both options are undervalued.
When the market prices converge to the prices with the realized volatility, both the call
and the put will become more valuable.

Although initially a straddle with an at-the-money call and an at-the-money put ($K=S$)
has a delta close to 0, as the stock price moves away from the strike price, the delta is no
longer close to 0 and the investor is exposed to stock price movements. So a straddle is
not a pure bet on stock volatility. For a pure bet on volatility, it is better to use volatility
swaps or variance swaps.\footnote[11]{For detailed discussion about volatility swaps, please refer to the paper \emph{More Than You Ever Wanted to Know about Volatility Swaps} by Kresimir Demeterfi, et al. The paper shows that a variance swap can be approximated by a portfolio of straddles with proper weights inversely proportional to $1/K^2$.} For example, a
variance swap pays $N\times(\sigma_r^2-K_{\mathrm{var}})$, where $N$ is the notional
value, $\sigma_r^2$ is the realized variance and $K_{\mathrm{var}}$ is the strike for the
variance.

\subsubsection*{Binary options}

What is the price of a binary (cash-or-nothing digital) European call option on a
non-dividend paying stock if the stock price follows a geometric Brownian motion? How
would you hedge a cash-or-nothing call option and what's the limitation of your hedging
strategy?

\solution
A cash-or-nothing call option with strike price $K$ pays \$1 if the asset price is above
the strike price at the maturity date, otherwise it pays nothing. The price of the option is
$c_B=e^{-r\tau}N(d_2)$ if the underlying asset is a non-dividend paying stock. As we have
discussed in the derivation of the Black-Scholes formula, $N(d_2)$ is the probability that
a vanilla call option finishes in the money under the risk-neutral measure. So its discounted
value is $e^{-r\tau}N(d_2)$.

Theoretically, a cash-or-nothing call option can be hedged using the standard delta hedging
strategy. Since

\[
\Delta=\frac{\partial c_B}{\partial S}=e^{-r\tau}N'(d_2)\frac{1}{S\sigma\sqrt{\tau}},
\]

a long position in a cash-or-nothing call option can be hedged by shorting
\[
e^{-r\tau}N'(d_2)\frac{1}{S\sigma\sqrt{\tau}}
\]
shares (and a risk-free money market
position). Such a hedge works well when the difference between $S$ and $K$ is large and
$\tau$ is not close to 0. But when the option is approaching maturity $T$ ($\tau\to0$)

% ZHOU-SOURCE-PAGE: 177 PRINTED: 161

and the stock price $S$ is close to $K$, $\Delta$ is extremely volatile\footnote[12]{$S\to K$
and $\tau\to0\Rightarrow\ln(S/K)\to0\Rightarrow d_2\to(r/\sigma+0.5\sigma)\sqrt{\tau}\to0
\Rightarrow\Delta\to\dfrac{e^{-r\tau}}{\sqrt{2\pi}\sigma S\sqrt{\tau}}\to\infty$.} and small changes in the stock price cause very large changes in $\Delta$. In these cases, it is
practically impossible to hedge a cash-or-nothing call option by delta hedging.

We can also approximate a digital option using a bull spread with two calls. If call options
are available for all strike prices and there are no transaction costs, we can long
$1/2\varepsilon$ call options with strike price $K-\varepsilon$ and short $1/2\varepsilon$
call options with strike price $K+\varepsilon$. The payoff of the bull spread is the same as
the digital call option if $S_T\leq K-\varepsilon$ (both have payoff 0) or
$S_T\geq K+\varepsilon$ (both have payoff \$1). When $K-\varepsilon<S_T<K+\varepsilon$,
their payoffs are different. Nevertheless, if we set $\varepsilon\to0$, such a strategy will
exactly replicate the digital call. So it provides another way of hedging a digital call option.
This hedging strategy suffers its own drawback. In practice, not all strike prices are traded
in the market. Even if all strike prices were traded in the market, the number of options
needed for hedging, $1/2\varepsilon$, will be large in order to keep $\varepsilon$ small.

\subsubsection*{Exchange options}

How would you price an exchange call option that pays $\max(S_{T,1}-S_{T,2},0)$ at
maturity. Assume that $S_1$ and $S_2$ are non-dividend paying stocks and both follow
geometric Brownian motions with correlation $\rho$.

\solution
The solution to this problem uses change of numeraire. Numeraire means a unit of
measurement. When we express the price of an asset, we usually use the local currency as
the numeraire. But for modeling purposes, it is often easier to use a different asset as the
numeraire. The only requirement for a numeraire is that it must always be positive.

The payoff of the exchange option depends on both $S_{T,1}$ (price of $S_1$ at the
maturity date $T$) and $S_{T,2}$ (price of $S_2$ at $T$), so it appears that we need two
geometric Brownian motions:

\[
dS_1=\mu_1S_1dt+\sigma_1S_1dW_{t,1}
\]

\[
dS_2=\mu_2S_2dt+\sigma_2S_2dW_{t,2}.
\]

Yet if we use $S_1$ as the numeraire, we can convert the problem to just one geometric
Brownian motion. The final payoff is
$\max(S_{T,1}-S_{T,2},0)=S_{T,1}\max\left(\frac{S_{T,2}}{S_{T,1}}-1,0\right)$. When
